\section{What further constant tuning can and cannot do}\label{sec:envelope}

\subsection{Exact optimum of the preceding nine-plan envelope}

This subsection uses the preceding packet's rectangular backend,
integer scale \(1400000\), and its nine plans. It helps distinguish a new
plan construction from merely reweighting the old certificate.
Let \(a_i\) be the chosen support with index
\[
 (67,67,69,69,67,69,67,69,69)_i,
\]
using the supplied lexicographic conventions. Put
\(G_{ij}=a_i\cdot a_j\), \(z=G^{-1}\mathbf1\), and
\(p=z/(\mathbf1^{\mathsf T}z)\).
The accompanying rational checker proves that \(G\) is invertible and all
\(p_i>0\). Define
\[
 H=\frac{1}{1400000^2\,\mathbf1^{\mathsf T}G^{-1}\mathbf1}.
\]
With these weights, conditional pair elimination leaves 1536 support
tuples. Their exact maximum weighted squared norm is \(H\).

\begin{proposition}[Restricted nine-plan optimum]\label{prop:nineopt}
For the old nine menus with their old backend, the worst-case value of
\(\min_i F_i(g)\) over nonnegative unit Euclidean vectors \(g\) is exactly
\(\sqrt H\). At the retained root schedule its coefficient is strictly
between \(29.746057\) and \(29.746058\).
\end{proposition}
\begin{proof}
The weighted version of Lemma~\ref{lem:norm}, verified globally with
the rational \(p\), gives the upper bound.
For the lower bound put
\[
 b=\sum_i p_i a_i/1400000.
\]
The equation \(Gp=\mathbf1/(\mathbf1^{\mathsf T}z)\) implies
\[
 \norm b^2=H,\qquad (a_i/1400000)\cdot b=H
 \quad\hbox{for every }i.
\]
All coordinates of \(b\) are nonnegative. At \(g=b/\sqrt H\),
each plan has a support of value \(\sqrt H\), so
\(\min_iF_i(g)\ge\sqrt H\).
Together with the upper bound this proves equality.
\end{proof}

This lower bound belongs to the Euclidean support envelope. It is not a
SAT lower bound and need not be attained by actual generated mode lists.
It nevertheless proves that further manipulation of mixing weights,
even depending on the support tuple, cannot improve that unchanged
nine-plan envelope below the displayed value. The new affine plans
escape the premise by changing the menu.

\subsection{A uniform-entropy limit for this twelve-plan architecture}

For any ten-group plan ending in blocks of sizes \(4,3,3\), uniform
unordered edge entropy gives final logarithmic dimensions in the ratio
\(12:12:9\). With the anchored backend \eqref{eq:rectcost}, the prescribed
cost envelope of the last operation is at least the linear branch \(12+21c\)
times the entropy of one edge.
Consequently its Euclidean-envelope coefficient obeys
\begin{equation}\label{eq:uniformlower}
 \beta\ge\frac{\sqrt2}{10\sqrt{45}}(12+21c).
\end{equation}
This holds for every menu of such plans, because uniform entropy makes
all final partitions identical in dimension.
At \(\eta_0=301/20000\), the resulting root coefficient is at least
\[
 \frac{\kappa}{\sqrt{301/20000}}
 \frac{\sqrt2(12+21c)}{10\sqrt{45}}
 =29.0807310194429\ldots.
\]
Thus this backend, root schedule, and final block architecture leave only
about \(0.0682\) absolute coefficient headroom below \(C_*\).
Changing just these menus cannot reach the unchanged tail branch
\(29.0657485\ldots\).
This is a useful stopping criterion for that particular search, not a
limitation on other contraction architectures or improved matrix backends.

\section{Attempts to change the logarithmic order}\label{sec:order}

The main square-root scale comes from the mode entropy
\(K(M,u)=\Theta(M/\sqrt u)\) together with admissible
\(u=O(\log s)\). A fixed finite improvement of a positive local coefficient
keeps that scale. The extra \(\log\log n\) in the dense SAT reduction is
the information needed by the current selected-incidence encoding.
It is not a polynomial arithmetic factor that disappears by using a
faster multiplication routine.

Several apparent order improvements were already explored in earlier
packets. Retuning the Hermite taper, imposing the ordinary-degree cap
alone, applying the same contractor directly to encoded mode indices,
and optimizing ordinary sparsification against the current density cost
do not supply the missing recursive interface \cite{barriers,precision}.
We did not repackage those conditional routes as new theorems.

Two additional elementary obstructions emerged from the present attempts.
Their hypotheses are deliberately narrow.

\subsection{Changing the prime requires a different packet geometry}

\begin{lemma}[Two good cyclic shifts]\label{lem:prime}
Let \(p\) be prime and \(R,C:\mathbb Z/p\mathbb Z\to A\) be row and
column group labels. A shift \(d\) is good if
\(R(t)=C(t+d)\) for every \(t\).
If two distinct shifts are good, both labelings are the same constant.
\end{lemma}
\begin{proof}
For good \(d\ne e\),
\(C(t+d)=R(t)=C(t+e)\).
Thus \(C(s)=C(s+e-d)\) for every \(s\).
The nonzero difference \(e-d\) generates the cyclic group of prime
order, so iterating this equality makes \(C\) constant everywhere.
Then \(R(t)=C(t+d)\) gives the same constant.
\end{proof}

For \(p=2\) there is one nonzero shift, so a nonconstant swapped labeling
can localize every nonzero residue. For an odd prime, two such residues
already force constancy. A nonconstant prime-cycle construction can retain
at most one of the \(p-1\) nonzero residues as good.
Thus the observation that a larger prime reduces the number of precision
digits does not by itself generalize the binary packet proof.

The precise finite-group extension identifies an escape: for a finite
abelian group \(G\) and a nonempty set \(D\) of good shifts, \(C\) need only
be constant on cosets of
\(\langle d-e:d,e\in D\rangle\).
Conversely that condition, with \(R(t)=C(t+d_0)\) for a fixed \(d_0\in D\),
makes every shift in \(D\) good. This follows by invariance under the
generators in one direction and under \(d-d_0\) in the other.
Affine vector spaces and other packet geometries are therefore not ruled
out by Lemma~\ref{lem:prime}.

\subsection{Coherent original-variable substitutions do not automatically split}

\begin{lemma}[Incidence coherence]\label{lem:coherence}
Give variables and factors group labels which agree at every
factor--variable incidence. Then every connected component of the
bipartite incidence graph has a single group label.
\end{lemma}
\begin{proof}
The labels at the endpoints of every edge are equal. Equality propagates
along a path, and any two vertices of a connected component are joined
by a path.
\end{proof}

\begin{corollary}[Shared linear-form version]
Let \(E=\bigoplus_{j=1}^r E_j\). Replace every occurrence of original
variable \(v\) by the same nonzero vector or linear form \(\ell_v\in E\).
Suppose factor \(f\) is assigned to group \(\sigma(f)\), and locality is
enforced by \(\ell_v\in E_{\sigma(f)}\) for every variable in its scope.
Then all factors in a connected incidence component lie in one group.
If its \(k\) variable images are also linearly independent, that group
has dimension at least \(k\).
\end{corollary}
\begin{proof}
If a variable met factors in two different groups, its image would belong
to their intersection, which is zero, contradicting nonzeroness.
The incidence lemma then gives one group per connected component.
Linear independence of \(k\) vectors in that group implies the dimension
bound.
\end{proof}

This rejects only a direct, fully coherent substitution that enforces
locality by structural containment. It permits rank-deficient or vanishing
images, factor-specific cancellations, measurements, correction modes,
and occurrence-dependent substitutions. The source's measurement
machinery falls outside the restrictive hypothesis. The corollary is not
a barrier to all degree-preserving encodings.

\subsection{Decision alone does not justify tiny parity-check budgets}

Reducing the precision of the final answer is not enough to remove the
linear term in the number of parity tests. The sparse-query degree bound
already contains a term proportional to the number of tests, even when
the requested output precision is small.
A useful scoped example is a formula with disjoint clauses: every
Boolean clause-violation vector is realizable. If a linear parity map
accepts exactly the zero violation vector, its kernel must be trivial,
so its rank equals the number of clauses.
This does not make disjoint clauses difficult to solve directly. It only
invalidates a general argument that decision automatically permits
sublinearly many exact parity filters.

We have not proved a new smaller-order deterministic 3-SAT algorithm.
No experiments in this archive are evidence for such a claim.

